Penelitian Tugas Akhir ini tidak terlepas dari dukungan, kerja sama, dan bantuan dari berbagai pihak.
Oleh karena itu, penulis menyampaikan rasa terima kasih yang mendalam kepada:

\begin{enumerate}
  \item Kedua orang tua, Bapak Mahanto dan Ibu Purwanti, Kakak Ririn, serta Almh. Nenek Parinem yang senantiasa memberikan motivasi dan doa selama proses perkuliahan hingga penyusunan Tugas Akhir ini.
  \item Bapak Muhammad Syifa'ul Mufid, S.Si., M.Si., D.Phil. selaku dosen pembimbing Tugas Akhir yang telah memberikan bimbingan dan arahan selama penyusunan Tugas Akhir ini.
  \item Bapak Dr. mont. Kistosil Fahim, S.Si, M.Si, Ibu Soleha, S.Si, M.Si, dan Drs. Lukman Hanafi, M.Sc, sebagai dosen penguji Tugas Akhir yang telah memberikan masukan dan saran guna meningkatkan kualitas Tugas Akhir ini.
  \item Ibu Prof. Dr. Dra. Mardlijah, MT selaku dosen wali yang telah membimbing dan memberikan dukungan akademik selama perkuliahan.
  \item Kepala Departemen Matematika ITS dan Kepala Program Studi Sarjana Matematika ITS beserta jajarannya.
  \item Seluruh dosen dan tenaga kependidikan di Departemen Matematika ITS yang telah membimbing, membagikan ilmu dan pengalaman, serta memberikan bantuan selama masa studi.
  \item Teman-teman seperjuangan mahasiswa yang telah memberikan dukungan dan motivasi selama perkuliahan.
\end{enumerate}

Penulis menyadari bahwa Tugas Akhir ini masih memiliki banyak kekurangan.
Kritik dan saran akan sangat diharapkan guna penyempurnaan penulisan di masa depan.
Akhir kata, penulis berharap Tugas Akhir ini dapat menebar kebermanfaatan bagi semua pihak.
